% Exemplo: estruturas de dados com forest (árvore binária de busca), TikZ (lista encadeada) e
% algpseudocodex (pseudocódigo da busca binária, com as palavras-chave em português).
% Compilar: latexmk -pdf estruturas-de-dados.tex  Testado com: forest 2.1, algpseudocodex 1.1, TeX Live 2024 ou mais recente
%%BEGIN:edpreambulo
\documentclass[12pt, a4paper]{article}
\usepackage[T1]{fontenc}
\usepackage[brazilian]{babel}
\usepackage{forest}                  % árvores
% nós com duas partes (lista)
\usetikzlibrary{shapes.multipart, positioning, calc}
\usepackage{algorithm}               % float "Algoritmo"
\usepackage{algpseudocodex}          % pseudocódigo
\floatname{algorithm}{Algoritmo}
% palavras-chave em português
\algrenewcommand\algorithmicfunction{\textbf{função}}
\algrenewcommand\algorithmicwhile{\textbf{enquanto}}
\algrenewcommand\algorithmicdo{\textbf{faça}}
\algrenewcommand\algorithmicif{\textbf{se}}
\algrenewcommand\algorithmicthen{\textbf{então}}
\algrenewcommand\algorithmicelse{\textbf{senão}}
\algrenewcommand\algorithmicreturn{\textbf{retorne}}
%%END:edpreambulo

\begin{document}
\pagestyle{empty}

%%BEGIN:edarvore
\begin{center}
\begin{forest}
  for tree={circle, draw, minimum size=2.2em, inner sep=1pt, s sep=1.5em}
  [50 [30 [20] [40]] [70 [60] [80, fill=yellow!40]]]
\end{forest}
\end{center}
%%END:edarvore

%%BEGIN:edlista
\begin{center}
\begin{tikzpicture}[
    no/.style={draw, rectangle split, rectangle split parts=2,
               rectangle split horizontal, minimum height=2.2em},
    seta/.style={-stealth, thick}]
  \node[no] (a) {12 \nodepart{two} \phantom{0}};
  \node[no, right=1cm of a] (b) {7 \nodepart{two} \phantom{0}};
  \node[no, right=1cm of b] (c) {31 \nodepart{two} \phantom{0}};
  \node[left=1cm of a] (ini) {\texttt{início}};
  \node[right=1cm of c] (nulo) {\textsc{nulo}};
  \draw[seta] (ini) -- (a);
  \foreach \x/\y in {a/b, b/c, c/nulo}
    \draw[seta] ($(\x.one split)!0.5!(\x.east)$) -- (\y);
\end{tikzpicture}
\end{center}
%%END:edlista

%%BEGIN:edalgoritmo
\begin{algorithm}[htbp]
  \caption{Busca binária em um vetor ordenado}
  \begin{algorithmic}[1]
    \Function{Busca}{$v, n, x$}
      \State $i \gets 1$; \ $f \gets n$
      \While{$i \le f$}
        \State $m \gets \lfloor (i + f)/2 \rfloor$
        \If{$v[m] = x$}
          \State \Return $m$
        \ElsIf{$v[m] < x$}
          \State $i \gets m + 1$
        \Else
          \State $f \gets m - 1$
        \EndIf
      \EndWhile
      \State \Return $-1$     \Comment{não encontrado}
    \EndFunction
  \end{algorithmic}
\end{algorithm}
%%END:edalgoritmo

\end{document}
